% Options for packages loaded elsewhere
\PassOptionsToPackage{unicode}{hyperref}
\PassOptionsToPackage{hyphens}{url}
\PassOptionsToPackage{dvipsnames,svgnames,x11names}{xcolor}
%
\documentclass[
  11pt,
  a4paper,
]{ltjsarticle}
\usepackage{amsmath,amssymb}
\usepackage{iftex}
\ifPDFTeX
  \usepackage[T1]{fontenc}
  \usepackage[utf8]{inputenc}
  \usepackage{textcomp} % provide euro and other symbols
\else % if luatex or xetex
  \usepackage{unicode-math} % this also loads fontspec
  \defaultfontfeatures{Scale=MatchLowercase}
  \defaultfontfeatures[\rmfamily]{Ligatures=TeX,Scale=1}
\fi
\usepackage{lmodern}
\ifPDFTeX\else
  % xetex/luatex font selection
\fi
% Use upquote if available, for straight quotes in verbatim environments
\IfFileExists{upquote.sty}{\usepackage{upquote}}{}
\IfFileExists{microtype.sty}{% use microtype if available
  \usepackage[]{microtype}
  \UseMicrotypeSet[protrusion]{basicmath} % disable protrusion for tt fonts
}{}
\makeatletter
\@ifundefined{KOMAClassName}{% if non-KOMA class
  \IfFileExists{parskip.sty}{%
    \usepackage{parskip}
  }{% else
    \setlength{\parindent}{0pt}
    \setlength{\parskip}{6pt plus 2pt minus 1pt}}
}{% if KOMA class
  \KOMAoptions{parskip=half}}
\makeatother
\usepackage{xcolor}
\usepackage[margin=24mm]{geometry}
\setlength{\emergencystretch}{3em} % prevent overfull lines
\providecommand{\tightlist}{%
  \setlength{\itemsep}{0pt}\setlength{\parskip}{0pt}}
\setcounter{secnumdepth}{5}
\ifLuaTeX
  \usepackage{selnolig}  % disable illegal ligatures
\fi
\IfFileExists{bookmark.sty}{\usepackage{bookmark}}{\usepackage{hyperref}}
\IfFileExists{xurl.sty}{\usepackage{xurl}}{} % add URL line breaks if available
\urlstyle{same}
\hypersetup{
  colorlinks=true,
  linkcolor={blue},
  filecolor={Maroon},
  citecolor={Blue},
  urlcolor={blue},
  pdfcreator={LaTeX via pandoc}}

\author{}
\date{}

\begin{document}

\hypertarget{ux4e00ux6b21ux72ecux7acb}{%
\section{一次独立}\label{ux4e00ux6b21ux72ecux7acb}}

ベクトル \texttt{v\_1,...,v\_k} が一次独立とは

\[
a_1v_1+\cdots+a_kv_k=0
\]

を満たす係数が

\[
a_1=\cdots=a_k=0
\]

しかないことです。

\hypertarget{ux76f4ux611f}{%
\subsection{直感}\label{ux76f4ux611f}}

一次独立とは、「どのベクトルも他のベクトルの組合せで作れない」ことです。つまり各ベクトルが新しい方向の情報を持っています。

一次従属なら、少なくとも一つは冗長です。

\hypertarget{ux884cux5217ux3068ux306eux95a2ux4fc2}{%
\subsection{行列との関係}\label{ux884cux5217ux3068ux306eux95a2ux4fc2}}

\texttt{A={[}v\_1\ ...\ v\_k{]}} とすれば、一次独立性は

\[
Ax=0
\]

が零解だけを持つことと同値です。

この事実は後に

\begin{itemize}
\tightlist
\item
  full column rank
\item
  逆行列
\item
  固有ベクトル
\item
  SVD の零特異値
\end{itemize}

と接続します。

\hypertarget{ux3088ux304fux3042ux308bux8aa4ux89e3}{%
\subsection{よくある誤解}\label{ux3088ux304fux3042ux308bux8aa4ux89e3}}

直交していれば一次独立ですが、一次独立だから直交しているとは限りません。独立性は代数的概念、直交性は内積を導入した後の幾何学的概念です。

\end{document}
